\documentclass[11pt,a4paper]{article}
\usepackage[left=0.75in,top=0.75in,right=0.75in,bottom=0.75in]{geometry}
\usepackage{enumitem}
\usepackage{hyperref}
\usepackage{titlesec}
\usepackage{xcolor}

\definecolor{primary}{RGB}{0, 51, 102}
\definecolor{secondary}{RGB}{90, 90, 90}

\titleformat{\section}{\Large\bfseries\color{primary}}{}{0em}{}[\titlerule]
\setlist[itemize]{noitemsep, topsep=2pt, leftmargin=*}

\newcommand{\entry}[4]{
    \noindent \textbf{#1} \hfill {\color{secondary}#2} \\
    \textit{#3} \hfill \textit{\color{secondary}#4}\vspace{1mm}
}

\begin{document}

%======================
% Header
%======================
\noindent {\Huge \textbf{\color{primary} Ramazan Bora}} \\[4pt]
Brisbane, QLD, Australia \hfill +61 421 099 850 \\
\href{mailto:r.bora@uq.edu.au}{r.bora@uq.edu.au} \hfill \href{https://ramazanbora.com}{ramazanbora.com} \\[-2pt]

\vspace{3mm}

%======================
% Profile
%======================
\section{Profile}

Economics PhD candidate in empirical economics and political economy, with a research agenda centred on using granular open data across space, time, and behaviour to study institutions, governance, and decision-making. My work combines original data construction from fragmented administrative, geospatial, remote-sensing, and behavioural sources with rigorous quantitative methods, including difference-in-differences, event studies, and high-dimensional fixed effects. Experienced in delivering end-to-end empirical research pipelines and policy-facing data products for academic and government stakeholders.

%======================
% Research interests & methods
%======================
\section{Research Interests \& Methods}
\textbf{Fields:} Political economy; institutions and governance; applied economics; computational social science. \\
\textbf{Methods:} Causal inference; panel data econometrics; GIS; remote sensing \& night lights; administrative data; large-scale behavioural data; algorithmic benchmarking; open-source data.

%======================
% Education
%======================
\section{Education}

\entry{University of Queensland}{Brisbane, Australia}{Ph.D. in Economics (thesis under examination; expected completion 2026)}{2021 -- present}
\begin{itemize}
    \item \textbf{Thesis:} \textit{Granularity in time, space, and behavior: Essays in empirical economics using open data}.
\end{itemize}

\entry{Penn State University \& Ibn Haldun University}{USA \& Turkey}{Doctoral studies in Economics}{2010 -- 2021}

\entry{Ko\c{c} University}{Istanbul, Turkey}{M.A. in Economics}{2008 -- 2010}

\entry{Ko\c{c} University}{Istanbul, Turkey}{B.Sc. in Electrical \& Electronics Engineering}{2003 -- 2008}

%======================
% Research
%======================
\section{Research}


\subsection*{Job Market Paper}

\begin{itemize}
    \item \textbf{A High-Resolution View of the Hinterland: Local Government Centralization and Rural Development in Turkey.} \\
    Examines the causal effects of Turkey’s 2012 metropolitan municipality reform on rural development. Reconstructs high-resolution village boundary maps and links them to VIIRS night lights, population outcomes, service-delivery proxies, and driving-distance measures. Implements a contiguous-border difference-in-differences design to estimate development effects, population retention, and distance-related penalties in service reach.
\end{itemize}

\subsection*{Working Papers}

\begin{itemize}
    \item \textbf{Birthplace Favoritism Revisited: Replication, Modern Evidence, and Crisis Dynamics.} \\
    \textit{UQ School of Economics Discussion Paper Series, 2025; working paper.} \\
    Replicates and extends classic evidence on leader birthplace favoritism using open-source data, VIIRS night lights, and high-dimensional fixed effects. Exploits monthly temporal resolution to study leadership transitions and crisis dynamics, showing that favoritism emerges immediately upon accession and intensifies during the COVID-19 period, particularly in weaker democracies.
\end{itemize}

\subsection*{Works in Progress}

\begin{itemize}
    \item \textbf{Measuring Cognition at Scale: The Association Between Ramadan Exposure and Chess Performance.} \\
    Builds a move-level panel of approximately 16 million moves from 800,000 online chess games, benchmarking decisions with the Stockfish engine. Estimates changes in move quality, speed, and play timing using difference-in-differences with player and period fixed effects, documenting improved move quality, faster decisions, and systematic shifts toward nocturnal play during Ramadan exposure.

    \item \textbf{A Novel Integerisation Procedure for Population Data} (with P. Wohland-Jakhar). \\
    Develops an integerisation method for IPF-based small-area estimation that reconciles integer and marginal constraints in sparse geographies, motivated by Indigenous population estimation in Queensland. Accompanied by reproducible implementation and stakeholder-facing audit tools.
\end{itemize}

%======================
% Conference Presentations
%======================
\section{Conference Presentations}

\begin{itemize}
    \item \textbf{38th Annual PhD Conference in Economics and Business} \hfill November 2024, ANU \\
    Presented: \textit{A High-Resolution View of the Hinterland: Local Government Centralization and Rural Development in Turkey}.

    \item \textbf{2025 Conference of the Society for Institutional \& Organizational Economics (SIOE)} \hfill August 2025, UNSW \\
    Presented: \textit{A High-Resolution View of the Hinterland: Local Government Centralization and Rural Development in Turkey}.

    \item \textbf{UQ Applied Micro Day} \hfill September 2025, UQ \\
    Presented: \textit{Birthplace Favoritism Revisited: Replication, Modern Evidence, and Crisis Dynamics}.
\end{itemize}

%======================
% Experience
%======================
\section{Research \& Applied Data Experience}

\textbf{Queensland Centre for Population Research (School of the Environment, UQ)} \hfill Brisbane \\
\textit{Senior Research Assistant} \hfill 2022 -- 2025
\begin{itemize}
    \item Developed a reproducible demographic estimation pipeline for the Queensland Government Statistical Office (QGSO), producing small-area (SA2) population estimates with a focus on Indigenous populations.
    \item Implemented a modified IPF procedure for non-full-dimensional constraint structures.
    \item Designed a novel integerisation algorithm converting fractional IPF outputs into integer population counts.
    \item Liaised with QGSO stakeholders and authored technical documentation and handover materials to support maintenance, auditability, and reuse.
\end{itemize}

\textbf{UQ International Comparisons Database (UQICD)} \hfill Brisbane \\
\textit{Senior Research Assistant} \hfill 2022 -- 2025
\begin{itemize}
    \item Refactored and documented the core \textbf{R} codebase supporting UQICD indicators and workflows (\href{https://uqicd.economics.uq.edu.au}{uqicd.economics.uq.edu.au}), improving maintainability and reproducibility.
    \item Built data pipelines enabling user-defined regional aggregations for purchasing power parity comparisons and developed publicly available dashboards in \textbf{R}.
    \item Managed acquisition, cleaning, harmonisation, and validation of heterogeneous international datasets to support stable updates.
\end{itemize}

\textbf{Real Estate \& Transport Project (Prof. Alicia Rambaldi)} \hfill Brisbane \\
\textit{Senior Research Assistant} \hfill 2025 -- present
\begin{itemize}
    \item Acquired and cleaned large administrative property transaction data in collaboration with the NSW Valuer General.
    \item Integrated sales records with GTFS transport networks and GIS layers, constructing accessibility measures for econometric modelling.
    \item Produced distance-weighted and route-diversity indices to quantify public transport access in commercial property analysis.
\end{itemize}

\textbf{Ibn Haldun University} \hfill Istanbul, Turkey \\
\textit{Research Assistant} \hfill 2018 -- 2021
\begin{itemize}
    \item Built text analytics and sentiment analysis workflows for social media data (Twitter API), including data collection, cleaning, and visualisation.
    \item Supported faculty research projects.
\end{itemize}

\textbf{Penn State University} \hfill State College, USA \\
\textit{Research Assistant (Prof. Bill Goffe; Prof. Vijay Krishna)} \hfill 2010 -- 2017
\begin{itemize}
    \item Developed a Python tool to flag irregularities in assessment data, automating detection of potential academic misconduct.
    \item Contributed proofs for research on voting rules and welfare properties.
\end{itemize}

\textbf{Ko\c{c} University (Engineering)} \hfill Istanbul, Turkey \\
\textit{Research Assistant} \hfill 2004 -- 2008
\begin{itemize}
    \item Applied pattern recognition methods to sensor data for driver identification (DRIVE-SAFE project).
    \item Implemented video deinterlacing algorithms for broadcast enhancement and signal-processing applications.
\end{itemize}

%======================
% Teaching
%======================
\section{Teaching Experience (Selected)}

\entry{University of Queensland}{Brisbane}{Senior Tutor / Tutor}{2021 -- present}
\begin{itemize}
    \item Senior Tutor: Mathematical Economics; Tools of Economic Analysis (tutorial coordination, administration, liaison with lecturers).
    \item Tutor: Advanced Microeconomics (M.A.); Strategy (M.A.); Managerial Economics; Mathematical Economics.
\end{itemize}

\entry{Ibn Haldun University}{Istanbul}{Teaching Assistant}{2018 -- 2021}
\begin{itemize}
    \item Tutor: PhD-level Microeconomics and Macroeconomics.
\end{itemize}

\entry{Penn State University}{USA}{Instructor of Record, Teaching Assistant, Tutor}{2010 -- 2017}
\begin{itemize}
    \item Instructor: Intermediate Microeconomics; Introductory Econometrics.
    \item Teaching Assistant: Microeconomics; Game Theory; Money and Banking; Economics of Law and Organization.
    \item Tutor, Morgan Center: One-on-one support for student athletes across economics, mathematics, and statistics.
\end{itemize}

%======================
% Skills
%======================
\section{Technical Skills}
\begin{itemize}
    \item \textbf{Programming:} R, Python, Stata, MATLAB; working knowledge of SQL.
    \item \textbf{Methods \& Data:} Causal inference; panel data analysis; administrative data; GIS; remote sensing; data scraping; algorithmic benchmarking.
    \item \textbf{Tools:} Google Earth Engine, AWS, \LaTeX.
\end{itemize}

%======================
% References
%======================
\section{References}
\begin{itemize}
    \item \textbf{Prof. Alicia Rambaldi} (Supervisor), University of Queensland. %\href{mailto:a.rambaldi@uq.edu.au}{a.rambaldi@uq.edu.au}
    \item \textbf{Dr. Haishan Yuan} (Supervisor), University of Queensland. %\href{mailto:h.yuan@uq.edu.au}{h.yuan@uq.edu.au}
    \item \textbf{Dr. Pia Wohland} (Project Lead, QGSO), University of Queensland. %\href{mailto:pia.wohland@uq.edu.au}{pia.wohland@uq.edu.au}
\end{itemize}

\end{document}